% Prose rendering of PrimeTensor/Analysis/Power/Growth.lean.
% Fragment: no \documentclass, no preamble, no document environment.

\section{End-to-end response growth for positive-depth powers}
\label{Analysis:Power:Growth:sec}

\leanfile{PrimeTensor/Analysis/Power/Growth.lean}

\subsection*{What this file establishes}

\texttt{PrimeTensor/Analysis/Growth.lean} replaced the chain rule's transport
assumption by a first-order growth condition, and verified that condition for
the identity and for constants --- the two maps where it is vacuous.  This file
supplies the first case with content: \emph{every positive-depth power is
response-scale controlled}.  With it, all three hypotheses the chain rule
places on an inner function are discharged for $y \mapsto y^{d}$, and the file
closes by stating the resulting rule.

A power costs refinement.  The $d$-th power is $\size{d} - 1$ multiplications,
each of which spends one level by
\texttt{PrimeTensor/Analysis/Near/Laws.lean}, so controlling $h^{d}$ at level
$\ell$ requires controlling $h$ at level $\ell$ refined $\size{d}-1$ times.
Recovering the level afterwards is a \emph{subtraction}, and $\name{Depth}$ has
none.  The first half of the file builds the two transforms that encode the
cost without one, and
Design note~\ref{Analysis:Power:Growth:note:adjoint} is about the identity that
makes them fit together.

\subsection{Encoding a finite refinement cost}

\begin{definition}[\lean{Depth.powerInput} and \lean{Depth.powerOutput}]
\label{Analysis:Power:Growth:def:transforms}
\begin{align*}
  \name{powerInput}(\ctor{one}, \ell) &= \ell, &
  \name{powerInput}(\ctor{succ}\,d, \ell)
    &= \name{powerInput}(d, \ctor{succ}\,\ell), \\
  \name{powerOutput}(\ctor{one}, \ell) &= \ell, &
  \name{powerOutput}(\ctor{succ}\,d, \ctor{one}) &= \ctor{one}, \\
  &&\name{powerOutput}(\ctor{succ}\,d, \ctor{succ}\,\ell)
    &= \name{powerOutput}(d, \ell).
\end{align*}
\end{definition}

\begin{leancode}
def powerInput : Depth → Depth → Depth
  | .one, level => level
  | .succ d, level => powerInput d (.succ level)
\end{leancode}

Numerically,
\[
  \size{\name{powerInput}(d,\ell)} = \size{\ell} + \size{d} - 1,
  \qquad
  \size{\name{powerOutput}(d,\ell)}
    = \max\bigl(\size{\ell} - (\size{d}-1),\; 1\bigr).
\]
The first adds the power's cost to a requested level; the second removes it,
saturating at the coarsest level rather than running off the bottom, since as
the docstring says there is no scale below $\ctor{one}$.  Both are defined by
recursion and neither mentions a difference: $\name{powerInput}$ moves a
$\ctor{succ}$ from the power onto the level, and $\name{powerOutput}$ cancels
$\ctor{succ}$ against $\ctor{succ}$ and stops when the level runs out.

\begin{lemma}[\lean{Depth.powerInput\_succ} and \lean{powerInput\_atOrAfter}]
\label{Analysis:Power:Growth:lem:input}
$\name{powerInput}(d, \ctor{succ}\,\ell) =
 \ctor{succ}\,\name{powerInput}(d,\ell)$, and
$\name{AtOrAfter}\bigl(\ell,\, \name{powerInput}(d,\ell)\bigr)$.
\end{lemma}

\begin{proof}
Both by structural recursion on $d$.  The first says the two additions
commute; the second says adding a cost only refines, and follows from the
first by transitivity, one $\ctor{later}$ at a time.
\end{proof}

\begin{theorem}[\lean{Depth.powerInput\_output\_join}]
\label{Analysis:Power:Growth:thm:adjoint}
For all $d$ and $\ell$,
\[
  \name{powerInput}\bigl(d,\; \name{powerOutput}(d,\ell)\bigr)
    = \name{join}\bigl(\ell,\; \name{powerInput}(d, \ctor{one})\bigr).
\]
\end{theorem}

\begin{proof}
By simultaneous recursion on $d$ and $\ell$; three of the four cases are
$\ctor{rfl}$.  In the remaining case both arguments are successors, and the
chain is: unfold both transforms one step; move the exposed $\ctor{succ}$ out
of $\name{powerInput}$ by Lemma~\ref{Analysis:Power:Growth:lem:input}; apply
the recursive case; observe that $\ctor{succ}$ of a join is the join of the
successors; and move the $\ctor{succ}$ back inside $\name{powerInput}$.
\end{proof}

\begin{designnote}[Truncated subtraction, and why the floor is harmless]
\label{Analysis:Power:Growth:note:adjoint}
Numerically the theorem reads
\[
  \bigl(\size{\ell} - (\size{d}-1)\bigr)^{+} + (\size{d}-1)
    = \max\bigl(\size{\ell},\, \size{d}\bigr),
\]
the standard identity for truncated subtraction: adding back what was
truncated recovers the original value, or the truncation threshold, whichever
is larger.  $\name{powerOutput}$ is a left adjoint to $\name{powerInput}$ in
the order $\name{AtOrAfter}$, and this is the composite of the adjunction on
one side.

The identity is exactly what the main theorem needs, and the reason is worth
stating.  Asked to control the response at level $\ell$, the proof offers the
response level $\name{powerOutput}(d,\ell)$.  For that to be justified the
perturbation must be controlled at $\name{powerInput}(d, \name{powerOutput}(d,
\ell))$, which by the theorem is not $\ell$ --- it is $\ell$ joined with the
power's minimum budget $\name{powerInput}(d,\ctor{one})$, that is, with
$\size{d}$.  When $\ell$ is already at least that fine the join is $\ell$ and
nothing is lost.  When it is not, the floor has bitten and more is being asked
of the perturbation than the caller supplied.

That is precisely why
Theorem~\ref{Analysis:Power:Growth:thm:main} opens by extracting an anchor
beyond which the perturbation is near the pivot at the minimum budget.  The
convergence hypothesis on $h$ pays for the floor, once, at the start; after
that the join is available for free at every stage and the adjunction closes.
\end{designnote}

\begin{lemma}[Three facts about \lean{Depth.advance}]
\label{Analysis:Power:Growth:lem:advance}
Advancing preserves the order in the base, and
\begin{align*}
  \name{advance}(\ctor{succ}\,\ell,\, g)
    &= \ctor{succ}\,\name{advance}(\ell,\,g), \\
  \name{advance}\bigl(\name{powerInput}(d,\ell),\, g\bigr)
    &= \name{advance}\bigl(\ell,\, \name{powerInput}(d,g)\bigr).
\end{align*}
The three are $\name{advance\_base\_atOrAfter}$,
$\name{advance\_succ\_base}$ and $\name{advance\_powerInput\_comm}$.
\end{lemma}

\begin{proof}
The first two are recursions on the gain.  The third is a recursion on $d$
using the first two; numerically both sides are
$\size{\ell} + \size{g} + \size{d} - 1$, so it is the commutativity and
associativity of adding three costs, proved by moving one $\ctor{succ}$ at a
time.
\end{proof}

\texttt{PrimeTensor/Analysis/Control.lean} proved that advancing respects the
order in the \emph{gain}; the first of these is the companion statement in the
base, and both are needed because the main proof advances a comparison between
two levels rather than between two gains.

\begin{lemma}[\lean{Depth.join\_eq\_left\_or\_right} and
  \lean{MulReal.scaleNear\_join}]\label{Analysis:Power:Growth:lem:join}
$\name{join}(a,b)$ is equal to $a$ or to $b$.  Consequently, if $x$ and $y$ are
near at level $a$ and also near at level $b$, they are near at
$\name{join}(a,b)$.
\end{lemma}

\begin{proof}
The first is a recursion on both arguments; being a maximum, the join is
whichever argument attains it, and the recursion produces the disjunction.  The
second then needs no combination argument at all: rewrite along whichever
disjunct holds and hand back the corresponding hypothesis.
\end{proof}

\begin{remark}[A third property of \lean{join}]\label{Analysis:Power:Growth:rem:join}
\texttt{PrimeTensor/Completion.lean} proved $\name{join}$ to be an upper bound
of its two arguments and remarked that nothing more was proved or needed.  Here
a third property appears, and it is stronger than it looks: a bound equal to
one of the things it bounds is a \emph{least} upper bound.  The file does not
say so, and uses the disjunction only to avoid having to combine two nearness
facts --- which, given that $\name{ScaleNear}$ has no transitivity, would
otherwise not be possible at all.
\end{remark}

\subsection{Powers on the carrier}

\begin{lemma}[\lean{MulReal.depthPow\_mul}]\label{Analysis:Power:Growth:lem:pow-mul}
$(ab)^{d} = a^{d} b^{d}$ for every depth $d$.
\end{lemma}

\begin{proof}
By recursion on $d$.  At $\ctor{one}$ both sides are $ab$.  At
$\ctor{succ}\,d$ the claim unfolds to
$(ab)(ab)^{d} = (a\,a^{d})(b\,b^{d})$; the recursive case rewrites the left
factor's tail, and the four-factor regrouping of
\texttt{PrimeTensor/Analysis/Rules.lean} finishes.
\end{proof}

\texttt{PrimeTensor/Analysis/Examples.lean} noted that $\name{depthPow}$ had
been given no laws.  This is the first of them, and the only one the file
needs; $a^{m}a^{n} = a^{m+n}$ is still absent, and still unstatable, since
$\name{Depth}$ has no addition.

\begin{theorem}[\lean{MulReal.scaleNear\_depthPow}]
\label{Analysis:Power:Growth:thm:pow-near}
If $h$ is near the pivot at $\name{powerInput}(d,\ell)$, then $h^{d}$ is near
the pivot at $\ell$.
\end{theorem}

\begin{proof}
By recursion on $d$.  At $\ctor{one}$ the power is $h$ and the input level is
$\ell$, so the hypothesis is the conclusion.

At $\ctor{succ}\,d$ the hypothesis is nearness at
$\name{powerInput}(d, \ctor{succ}\,\ell)$.  Weakening it by
Lemma~\ref{Analysis:Power:Growth:lem:input} and
\texttt{PrimeTensor/Analysis/Control.lean} gives $h$ near the pivot at
$\ctor{succ}\,\ell$; and the recursive case, applied to the same hypothesis at
level $\ctor{succ}\,\ell$, gives $h^{d}$ near the pivot at $\ctor{succ}\,\ell$.
Two facts at $\ctor{succ}\,\ell$ multiply to one at $\ell$, and
$h \cdot h^{d} = h^{\ctor{succ}\,d}$.
\end{proof}

This is the quantitative statement that a power costs refinement, and the
budget is exact: $\size{d}-1$ multiplications, one level each, is precisely the
$\size{d}-1$ that $\name{powerInput}$ adds.

\subsection{Response-scale control of a power}

\begin{lemma}[\lean{response\_depthPow}]\label{Analysis:Power:Growth:lem:response}
$\name{response}(y \mapsto y^{d},\, x,\, h)_n = (h_n)^{d}$.
\end{lemma}

\begin{proof}
The response is $\name{ratio}\bigl((x h_n)^{d},\, x^{d}\bigr)$; expand the
numerator by Lemma~\ref{Analysis:Power:Growth:lem:pow-mul} into
$x^{d}\,(h_n)^{d}$, and the reconstruction identity of
\texttt{PrimeTensor/Analysis/Derivative.lean} cancels $x^{d}$.
\end{proof}

An identity again, as
\texttt{PrimeTensor/Analysis/Examples.lean} predicted for the whole monomial
family.

\begin{theorem}[\lean{responseScaleControlledAt\_depthPow}]
\label{Analysis:Power:Growth:thm:main}
For every $x$ and every depth $d$, the map $y \mapsto y^{d}$ is response-scale
controlled at $x$.
\end{theorem}

\begin{proof}
Let $h$ approach the pivot and let a target gain $t$ be requested.  Ask the
convergence of $h$ for the level $\name{powerInput}(d,\ctor{one})$ --- the
power's minimum budget --- and take the anchor it returns.  Offer
$\name{powerInput}(d,t)$ as the response gain.

Beyond the anchor, suppose $h_n$ is near the pivot at level $\ell$.  It is also
near the pivot at the minimum budget, so by
Lemma~\ref{Analysis:Power:Growth:lem:join} it is near at
$\name{join}(\ell, \name{powerInput}(d,\ctor{one}))$.  Offer
$\rho = \name{powerOutput}(d,\ell)$ as the response level.  Two things must be
checked.

\emph{The response is near the pivot at $\rho$.}  By
Theorem~\ref{Analysis:Power:Growth:thm:adjoint} the joined level just obtained
\emph{is} $\name{powerInput}(d,\rho)$, so
Theorem~\ref{Analysis:Power:Growth:thm:pow-near} puts $(h_n)^{d}$ near the
pivot at $\rho$, and Lemma~\ref{Analysis:Power:Growth:lem:response} identifies
that with the response.

\emph{The comparison holds.}  Again by
Theorem~\ref{Analysis:Power:Growth:thm:adjoint},
$\name{powerInput}(d,\rho)$ is a join with $\ell$ as one argument, so $\ell$
lies at or before it.  Advance both by $t$, using the base-monotonicity of
Lemma~\ref{Analysis:Power:Growth:lem:advance}, and then commute the advance
past the input transform by the third part of the same lemma:
\[
  \name{advance}\bigl(\name{powerInput}(d,\rho),\, t\bigr)
  = \name{advance}\bigl(\rho,\, \name{powerInput}(d,t)\bigr),
\]
whose right-hand side is $\rho$ advanced by the offered response gain.  That
is exactly the required
$\name{AtOrAfter}(\name{advance}(\ell,t), \name{advance}(\rho, r))$.
\end{proof}

\begin{corollary}[\lean{hasMulDerivativeAt\_comp\_depthPow}]
\label{Analysis:Power:Growth:cor:chain}
If $g$ has derivative $D_g$ at $x^{d}$ and $D_g$ is scale controlled, then
$y \mapsto g(y^{d})$ has derivative
$\name{comp}\bigl(D_g,\, \name{depthPow}(d)\bigr)$ at $x$.
\end{corollary}

\begin{proof}
Apply the chain rule of \texttt{PrimeTensor/Analysis/Growth.lean}.  Its three
hypotheses on the inner map are now all available: the derivative from
\texttt{PrimeTensor/Analysis/Examples.lean}, the tangent control from
\texttt{PrimeTensor/Analysis/Control.lean}, and the response-scale control from
Theorem~\ref{Analysis:Power:Growth:thm:main}.
\end{proof}

\begin{remark}[The first chain rule with content]
\label{Analysis:Power:Growth:rem:first}
Both \texttt{PrimeTensor/Analysis/Chain.lean} and
\texttt{PrimeTensor/Analysis/Growth.lean} closed with the observation that the
rule could be discharged only for inner maps --- the identity, then also the
constants --- where it says nothing.  Corollary~\ref{Analysis:Power:Growth:cor:chain}
is the first instance that is not of that kind: the inner map genuinely
distorts scale, by $\size{d}-1$ levels, and the hypotheses are met anyway
because the distortion is finite and exactly accounted.  Nothing about $g$ is
assumed beyond differentiability and control of its tangent --- in particular
$g$ needs no response-scale control of its own, the rule requiring that only of
the inner factor.
\end{remark}

\begin{remark}[Still no iteration]\label{Analysis:Power:Growth:rem:no-iterate}
Applying the corollary twice would require
$\name{comp}\bigl(D_g, \name{depthPow}(d)\bigr)$ to be scale controlled, and
composition is still not shown to preserve that property --- the gap named in
\texttt{PrimeTensor/Analysis/Chain.lean} and untouched here.  Likewise
$\name{ResponseScaleControlledAt}$ is proved for powers, for the identity and
for constants, but for no product, inverse or composite of them, so a general
monomial $y \mapsto c\,y^{d}$ is not covered even though each of its factors
is.
\end{remark}

\subsection*{Summary of the correspondence}

\begin{center}
\small
\begin{tabular}{@{}lll@{}}
\hline
Lean declaration & Kind & Here \\
\hline
\lean{Depth.powerInput}                 & definition & Def.~\ref{Analysis:Power:Growth:def:transforms} \\
\lean{Depth.powerOutput}                & definition & Def.~\ref{Analysis:Power:Growth:def:transforms} \\
\lean{Depth.powerInput\_succ}           & theorem    & Lem.~\ref{Analysis:Power:Growth:lem:input} \\
\lean{Depth.powerInput\_atOrAfter}      & theorem    & Lem.~\ref{Analysis:Power:Growth:lem:input} \\
\lean{Depth.powerInput\_output\_join}   & theorem    & Thm.~\ref{Analysis:Power:Growth:thm:adjoint} \\
\lean{Depth.advance\_base\_atOrAfter}   & theorem    & Lem.~\ref{Analysis:Power:Growth:lem:advance} \\
\lean{Depth.advance\_succ\_base}        & theorem    & Lem.~\ref{Analysis:Power:Growth:lem:advance} \\
\lean{Depth.advance\_powerInput\_comm}  & theorem    & Lem.~\ref{Analysis:Power:Growth:lem:advance} \\
\lean{Depth.join\_eq\_left\_or\_right}  & theorem    & Lem.~\ref{Analysis:Power:Growth:lem:join} \\
\lean{MulReal.scaleNear\_join}          & theorem    & Lem.~\ref{Analysis:Power:Growth:lem:join} \\
\lean{MulReal.depthPow\_mul}            & theorem    & Lem.~\ref{Analysis:Power:Growth:lem:pow-mul} \\
\lean{MulReal.scaleNear\_depthPow}      & theorem    & Thm.~\ref{Analysis:Power:Growth:thm:pow-near} \\
\lean{response\_depthPow}               & theorem    & Lem.~\ref{Analysis:Power:Growth:lem:response} \\
\lean{responseScaleControlledAt\_depthPow} & theorem & Thm.~\ref{Analysis:Power:Growth:thm:main} \\
\lean{hasMulDerivativeAt\_comp\_depthPow}  & theorem & Cor.~\ref{Analysis:Power:Growth:cor:chain} \\
\hline
\end{tabular}
\end{center}

\noindent
Every declaration of \texttt{PrimeTensor/Analysis/Power/Growth.lean} appears
above.  The file contains no \lean{sorry}, no axiom, and no unproved named
proposition.  Design note~\ref{Analysis:Power:Growth:note:adjoint} and
Remarks~\ref{Analysis:Power:Growth:rem:join},
\ref{Analysis:Power:Growth:rem:first}
and~\ref{Analysis:Power:Growth:rem:no-iterate} are stated for the reader and
are not declarations of the file; the numerical readings of the two transforms
are glosses and are not formalised.
